\section{Approximate tagged permutations at logarithmic width}\label{sec:width}

Fix the finite field and the maximal flat dimension $k$ of
Theorem~\ref{thm:protocol}. All logarithms in this section are binary.
The tensor $\operatorname{Tag}(h,w)$ tests whether a list of $n=2^h$
tagged $w$-bit records is a permutation of a prescribed list, with each
tag giving the position of its payload in that list
(Definition~\ref{def:tag}). For $w=h$, the two lists together expose
$3nh$ bits. We show that, for every $r>0$, a factor of order $h^r$
above this port count suffices to approximate the test. All data and
tag indices stay exposed, and the size includes the catalyst ports.

\par\Needspace{5\baselineskip}
\begin{theorem}[Approximation at logarithmic payload width]\label{thm:cheaptag}
Under the conclusion of Theorem~\ref{thm:protocol},
for each $r>0$ there is a constant $C$ such that, for all sufficiently
large $h$, $\operatorname{Tag}(h,h)$ admits a trivalent approximation
of size $O(2^h h^{1+r})$ and gap at least
\[
               2^h\,2^{-C(\log(2h))^2}.
\]
In particular its gap is at least $2^h2^{-C'\sqrt h}$ for a fixed $C'$.
\end{theorem}

The proof compares two ways to change the payload width. With wide
payloads, pruning the routing identity saves order $nw$ in the price of
each discarded term, enough to cover the price of summing over tags and
selectors. Each surviving term uses tests with shorter addresses.
Compaction limits the number of tests installed, and amplification
makes their combined error small enough for the parent test. In the
other direction, stable radix sorting implements the width-$h$ test by
packing narrow records into wide ones. The two bounds together force
the overhead exponent at width $h$ to vanish.

\subsection{Compacting active records}

The amplification step corrects a sparse set of records. It has a
fixed array of possible correction slots, but the active positions depend
on the input. We connect the slots to a smaller array of targets in
advance. If at most half the targets are needed, the control can assign
every active slot to a distinct target along these fixed links. Bounded
degree keeps the number of physical alternatives per position constant.
The construction uses probabilistic expansion (see \cite{HLW}).

\begin{lemma}[Compaction with slack]\label{lem:slack-compaction}
Fix $C\ge1$. For $m\le CK$ slots and $K$ targets there is a bipartite
graph with degrees bounded in terms of $C$, on which any set of at most
$K/2$ active slots has an injective assignment to neighboring targets.
Given the graph as fixed advice, the assignment can be chosen
deterministically in $O_C(K)$ word operations on $O(\log(2K))$-bit words.
\end{lemma}

\begin{proof}
For bounded $K$ use the complete bipartite graph. For large $K$, fix a small $\Delta>0$ depending
on $C$. Keep $R=\lfloor K/4\rfloor$ targets in reserve. A bounded-choice
graph on the other targets can make every $\lceil\Delta K\rceil$ slots
meet every set of $\lceil K/5\rceil$ targets. Indeed, give each slot
$d$ random choices. The probability of missing a fixed pair of sets
is at most $(1-1/5)^{d\lceil\Delta K\rceil}$, and there are at most
$2^{(C+1)K}$ pairs. Choose $d$ large enough. Process the active
slots in order, taking the first free neighbor if there is one.
At least $K/5$ targets remain permanently free on this side. No
unmatched slot is adjacent to any of them, so fewer than
$\lceil\Delta K\rceil$ slots remain.

Independently give each slot eight choices in the reserve. For
$1\le l\le\Delta K$, the probability that some $l$ slots have at most
$4l$ distinct neighbors is at most
\[
   \binom{m}{l}\binom{R}{4l}(4l/R)^{8l}
          \le \bigl(C_1(l/K)^3\bigr)^l
\]
after increasing $C_1=C_1(C)$. Here the exponent is $8-4-1=3$. Take $\Delta$ small enough that
$4\Delta K\le R$ and $C_1\Delta^3<1/2$ for large $K$.
Splitting the sum at $l=\sqrt K$ shows its convergence to zero.
To bound the target degrees in both random constructions, partition
the slots in each choice layer into blocks the size of the relevant
target set. Use an independent random permutation within each padded block. Within each block, sampling without replacement bounds the probability of
landing entirely in a given subset by the corresponding with-replacement
probability. Thus both bounds still hold, and a target occurs at most once
per block per layer.

Among $l>0$ remaining slots in the reserve graph some target has
degree one on their incidences, since there are more than $4l$
neighbors and only $8l$ incidences. Assign its unique slot to it
and delete that slot with its incidences. A fixed-order queue of
degree-one targets and incidence lists maintain this operation in
linear time, skipping stale queue entries. The assigned targets are
distinct. The queue, the incidence scans and the greedy first pass fit the word
bounds of Definition~\ref{def:protocol}.
\end{proof}


\subsection{Amplification by sparse exact corrections}

To amplify the gap for a batch of tagged tests, we will evaluate a
small fraction of the tests exactly. We begin with the cost of one such
evaluation.

A single exact tagged test costs $O(nwh)$ for $w\ge h$. Indeed, an AKS
sorting network uses $O(nh)$ comparators to sort the second list by tag
\cite{AKS}. Comparing two tags and routing their payloads costs
$O(h+w)$ bit gates, including fan-out. Compare the sorted tags with the
prescribed positions and the sorted payloads with the first list.
The resulting Boolean verification gives an exact trivalent tensor
network over every field. A stable sort on $2^d$ records of $v$ bits
also has a network consisting of one $\operatorname{Tag}(d,v)$ call and
$O(2^d(d+v))$ additional size. Guess an input position for each proposed
output, check the permutation of the payloads with the call, and check
nondecreasing keys with increasing input position at ties. Each valid
input--output pair has exactly one witness.

\begin{lemma}[Amplification]\label{lem:tag-amplification}
Let $X$ approximate $f=\operatorname{Tag}(h,w)$ with exposed data, size
$S$, and gap $D$ (using a trivalent realization of $X$), where $w\ge h$ and $D\ge4k(\log h+3)$. Fix $C'>0$.
For all sufficiently large $h$ and every integer $s\ge1$ with
$\log s\le h^{C'}$, the tensor $f^{\otimes s}$ has a trivalent
approximation of size at most $C_{\rm amp}sS$ and gap at least
$sD/(2h)$. The size constant is absolute. The threshold for $h$ may
depend on $C'$ and $k$.
\end{lemma}

\begin{proof}
Truncate the tensor-power expansion by the number of errors. The
discarded terms contain enough errors to give the amplified gap, and the
retained terms require exact corrections at only a small fraction of
positions. Let $I$ be the through identity on the catalyst of $X$ and put
$F_0=f\otimes I$, $E=F_0-X$. In $(X+E)^{\otimes s}$ retain precisely
the terms with at most $p=\lfloor s/h\rfloor$ copies of $E$. Give this
sum the catalyst format of $s$ copies of $X$.

Fix a price $H=c\log\varphi$ on the center and the $2k$ unit catalyst
sites, induced by any spectral $k$-flat. Its scale satisfies $c\le2k$, by the star through the center.
Normalize spectral values of one copy by $\varphi(I)$. The normalized
values of $F_0,E,X$ are at most $1,z,1+z$, respectively, where
$z=2^{-D/c}$. Here ordinary sum is a restriction of direct sum, and
multiplication by a nonzero scalar preserves the price. Exact $X$ can
be powered directly. For finite $D$ we have $z\le1/h$. Weight each
omitted term with $j$ errors by $(hz)^{-j}$, and use multiplicativity
and the binomial theorem. The sum of the normalized values of all
omitted terms is at most
\[
                 (hz)^{p+1}(1+z+1/h)^s .
\]
This bounds the spectral value of the error relative to
$\varphi(I)^s$. Since $p+1\ge s/h$ and
$\log(1+z+1/h)\le (2\log e)/h<3/h$, the price discount is at least
\[
         \frac{s}{h}\bigl(D-c(\log h+3)\bigr)\ge \frac{sD}{2h}.
\]

For the size assertion, encode each position by one of three choices,
namely ordinary $X$, a selected correction $F_0$, or a selected
correction $-X$. Reject strings with more than $p$ selected positions. This gives
one encoding for each summand of the retained expansion after expanding
$E=F_0-X$. Put a switchable copy of $X$ at every position. A switched-off
module has every internal bit fixed to zero and local value one on
that assignment. This is implemented with a copied enable bit on its
bounded-arity tables and equality vertices if a wire requires one.
A position assigned $F_0$ bypasses its own catalyst bits by through
wires and routes only its $O(nw)$ data bits to an exact $f$ module.
Assign the sign locally at each selected $-X$. All unused fields and
inactive lanes have a unique zero padding.

If $p=0$, use $X^{\otimes s}$.
Otherwise, repeatedly halve the list length, rounding up, until at most
$4p$ positions remain. At every actual halving there are at most $p$
active positions and the destination length is at least $2p$.
The bounded-degree compaction graph for ratio two and its deterministic linear-word-time assignment algorithm (Lemma~\ref{lem:slack-compaction}) send the active positions injectively to neighbors.
Use the chosen switches for the whole data word of a position.
The sum of the list lengths is $O(s)$, hence data routing, including
its interface to the local $X$ modules, costs $O(snw)$. Install exact,
switchable $f$ modules at the final list, costing $O(pnwh)=O(snw)$.
All graph descriptions have $O(s)$ words of advice. The configuration
algorithm and the check on the number of selected positions use
$O(s)$ operations on $O(\log(2s))$-bit words. Verify the deterministic memory trace by sorting its records by address
and time with a Batcher network \cite{Batcher}. A scan in this order
checks each read against the initialization or the most recent write, as
in RAM verification for proof systems \cite{TinyRAM}. Chronology, unused records, and final configurations
have fixed conventions, so each legal choice has one history.
This control costs $O(s\,\mathrm{polylog}(2s))=O(snw)$ because the fixed
power of $\log(2s)$ is eventually bounded by $2^h$ in the stated
range. The implied size constants are absolute. Each exposed network
has $S\ge nw$, and all enable wires and catalyst bypasses cost $O(sS)$.
\end{proof}

\subsection{The width profile and stable radix packing}

For a polynomial payload width $w=\lceil h^a\rceil$, measure the size
as an overhead above $nw$, keeping the required gap essentially
proportional to the list length. Put $L(h)=\log(2h)$. For $a\ge1$,
define $r(a)$ as the infimum of the
real $r$ such that for some $C>0$ and all sufficiently large $h$ the
tensor at width $w=\lceil h^a\rceil$ admits exposed approximations of
size $O(nw h^r)$ and gap $n2^{-C L(h)^2}$ or greater. Write
$r_*=r(1)$ and $r_\infty=\inf_{a\ge1}r(a)$. The port count and the
exact construction give $0\le r(a)\le1$. Pinning matching padding
bits to zero makes a wide test serve a narrower payload. For the
reverse operation, split the wide payload into slices, copy the tags
to every slice, and multiply the predicates. This product enforces
the permutation condition on the common tags and every payload equality. The tag copies cost $O(nh)$ per slice.
For any fixed $a'\ge a$ the number of slices is polynomial in $h$.
Lemma~\ref{lem:substitution} loses at most $2k$ times that number
in absolute gap, absorbed by increasing $C$. Thus
\[
                  0\le r(a)-r(a')\le a'-a .
\]
The same argument permits constant factors in the width at fixed $a$.

In a radix sort, the digit length controls the cost of sorting a block,
whereas the pack width controls the cost of moving many records with
the same digit. Optimizing these two costs separately gives an upper
bound on $r_*$ in terms of the profile at larger widths.

\begin{proposition}[Packing wide records]\label{prop:tag-width-packing}
For every fixed $a>1$,
\[
 r_*\le 1-\frac1a+\max\{r(a)/a,r_\infty\}.
\]
\end{proposition}

\begin{proof}
To check the parent tagged relation, stably radix-sort a combined list
of sources and targets
with their keys and record-type flags and use a final scan to check for one
matching payload of each type at each tag. This list has $\Theta(n)$
records of $O(h)$ bits. Take $k_0=\lfloor h^{1/a}\rfloor$ digit bits per
pass, proceeding from least significant to most significant. For any
fixed $B>1$, take a dyadic pack size $P$ in $[h^{B-1},2h^{B-1})$.

Pad to blocks of $2^{2k_0}$ slots and stably sort each block on its
current digit, placing dummies last. Round each nonempty digit run up
to a multiple of $P$ and place its first real record at a multiple of
$P$. The original slots already accommodate the old dummies. Rounding and
alignment require only $O((2^{k_0}+1)P)$ additional slots per block. Scans compute run lengths, rounded
lengths, and destinations (including for old dummies) in $O(h)$ per slot. A second tagged transport implements the alignment. For instance, give each source its destination as a tag (transpose the perspective on which list is in prescribed order). Positions not yet assigned are filled with padding sources in increasing order; this last condition is enforced by carrying source positions as part of the payload and checking on the output side that those sources are in order. When a
tagged call is padded to a power of two, pin the payload on both sides
of each extra slot to zero and fix its output tag to the address of its
own input slot. The padded matching thus extends the original one
uniquely. All real records in a pack now have the same digit. Obtain the key of each nonempty pack
by a scan, stably sort the packs by that key with empty packs last, and
retain the order of real records within the packs. The result is the
stable sort on this digit among real records; scans ignore vacant slots and move forward across them. The list length grows each pass by a fraction
$O(2^{-k_0}P)$ and an additive $O(2^{2k_0}+2^{k_0}P)$ for the partial
block. Over $O(h/k_0)$ passes it stays $\Theta(n)$, with lengths fixed
before receiving the data.

There are a constant number of collections of block calls per pass,
each of common address length $l=2k_0+O(1)$, payload width $O(h)$ and
$\Theta(n)$ slots in total. Implement a block at the bound for $r(a)+e$,
$e>0$, with matching constant-factor width adjustments. Amplify its
$s=\Theta(n/2^l)$ copies by Lemma~\ref{lem:tag-amplification}, using
$\log s\le h=l^{O_a(1)}$. Each collection costs
$O(nh\,h^{(r(a)+e)/a})$ and has gap at least
$\Omega(n2^{-C L(l)^2}/l)$ for a constant $C$.
Each pass also uses one pack sort with $\Theta(n/P)$ slots, address
length $h-O(\log h)$ and width $O(hP)$. The bound at $r(B)+e$ costs
$O(nh\,h^{r(B)+e})$ and gives gap at least $n2^{-C_1 L(h)^2}$ after
enlarging $C_1$ to absorb the factor $P$. The scans and tag copying cost
$O(nh)$ per pass. Treat each collection of block calls as one product
gate, as allowed for a multilinear tensor contraction. Across all passes
there are polynomially many substituted gates, and the gaps dominate
their number. A further increase of $C_1$ makes the result admissible in
the width profile. Multiply the size bounds by $O(h^{1-1/a})$, let
$e\downarrow0$, and take the infimum over $B>1$.
\end{proof}

\subsection{Pruning the routing identity and counting shorter calls}

For a valid tag assignment $\pi$, denote by $M_\pi$ the payload
matching tensor with entries $\prod_i\mathbf{1}_{q_i=p_{\pi(i)}}$,
and by $\delta_\pi$ the basis tensor supported on that assignment to
the tag ports. Definition~\ref{def:protocol} supplies an exact routing
identity for $M_\pi\otimes I$, where $I$ denotes the protocol catalyst
identity. Tensor each routing identity with $\delta_\pi$ at the origin
and sum over permutations. The resulting tensor is exactly
\[
 \sum_{\pi\in S_n}\delta_\pi\otimes M_\pi\otimes I
       =\operatorname{Tag}(h,w)\otimes I.
\]
In a network this sum is implemented by checking that the exposed tags
form a permutation and using those tags in the control. Each valid tag
assignment activates its own routing identity with the original selector
multiplicities. This representation also places the error from pruning
different tag assignments on the same set of exposed ports.

Fix $\kappa>0$ and $A>1+\kappa$, and take $h^{1+\kappa}\le w\le h^A$. Retain the terms for which
the sum of $\lambda y$ over links joining different layers is at least
$2MN_c/w-n$. To price one discarded term, place its sites at their layer
coordinates divided by $M$ in the unit catalyst configuration on any
spectral $k$-flat. Local tensors and same-layer links cost zero.
Every cross-layer link costs at most $\lambda yw/M$ since its charge
dominates its Euclidean layer distance. The local factor $\delta_\pi$
has price zero. Thus a discarded term has
price at most $2N_c-nw/M$. Across all valid permutations and all
selectors there are at most $2^{O(nh\,\mathrm{polylog}\,h)}$ terms.
The cost of their sum is at most their maximum price plus
$O_k(nh\,\mathrm{polylog}\,h)=o(nw)$. Consequently the surviving
contraction approximates $\operatorname{Tag}(h,w)$ with exposed gap
$\Omega(nw)$ relative to the through price $2N_c$. Testing the pruning
condition adds only linear word work using the fixed rational charges.
Each surviving term satisfies
\begin{equation}\label{eq:tag-positive-budget}
            \sum_{t>0}(\lambda y+t^p)\le 11n .
\end{equation}


We next turn the surviving traffic bound into a recursion for network
size. Let $u:[1,\infty)\to(0,\infty)$ be a smooth, eventually constant
function with $u(1)>r_*$ and $0\le-u'<1$. We will prove $r(a)\le u(a)$
by induction on address length. The width-$h$ estimate starts the
induction near $a=1$. For wider payloads, the shorter calls must fit
inside the parent allowance $nwh^{u(a)}$.

Write $a=\log_h w$, $b=-u'(a)$ and $v=u(a)+ab$, and fix
$\epsilon>0$. When a child call uses fractions $T$ and $Y$ of
the parent's address length and payload width, its size allowance
\emph{per slot}, relative to the parent's, is asymptotic to
\[
                          T^v Y^{1-b}.
\]
Indeed the width exponent changes from $a$ to $\log_{Th}(Yw)$.
The calculation after \eqref{eq:tag-uniform-width} gives the uniform
estimate. Since $v-b=u(a)+(a-1)b>0$, we can choose $Y=T^p$ with $p>1$
close to one so that $v-pb>0$ uniformly. Even with order $n/Y$ installed
slots the relative cost is then $O(T^{v-pb})$, which tends to zero
as $T$ decreases. Equation~\eqref{eq:tag-positive-budget}
bounds the active slots in each summand. Compaction turns this into a
bound on the number of slots installed in the common network.

\begin{lemma}[Shorter calls]\label{lem:tag-shorter}
One choice of protocol parameters works for all sufficiently
large $h$ and $1+\kappa\le a\le A$. The pruned contraction is implemented
exactly relative to the following tagged calls, with fixed positions and
wiring for each format.
\begin{enumerate}
\item A fixed finite menu of dyadic pairs $(T,Y)$ has $0<T\le1/2$,
$0<Y\le1$. At each installed type the calls have common address length
$h'=Th+o(h)$ and width $w'=\lceil Yw\rceil$. Counting all physical
slots, including idle slots, their numbers $S_{T,Y}$ satisfy
\[
       S_{T,Y}=\Omega(n)\quad\text{per type},\qquad
       \sum_{T,Y} S_{T,Y} T^v Y^{1-b}\le\epsilon n .
\]
\item A bounded number of control calls have address length
$h+O(\log h)$, width equal to address length, and
$O(n\,\mathrm{polylog}\,h)$ slots altogether.
\end{enumerate}
All other gates and wires, including the incidences at the calls, cost
$O(nw\,\mathrm{polylog}\,h)$.
\end{lemma}

\begin{proof}
The continuous function $v-b=u(a)+(a-1)b$ has positive minimum
on $[1+\kappa,A]$. Choose a rational $p>1$ sufficiently close to one
that $v-pb$ still has positive minimum. Take $W$ large enough for
Theorem~\ref{thm:protocol}, with $pW$ integral, and use a single type
$(T,Y)$ given by
\[
 T=2^{-W},\qquad Y=T^p.
\]
Scale-zero groups have bounded size and are wired with switches. For the scale-$T$ groups, orient links and reserve distinct source and destination slots. Split each nonempty link into transports of width $\lceil Y w\rceil$, padding its last chunk; protocol bounds allow us to reserve fixed numbers of slots for these chunks. By \eqref{eq:tag-positive-budget} at most $c_*n$ transports are used in any one summand, where $c_*=22 T^{-p}$. Hence for $W$ sufficiently large
\[
 7c_* T^v Y^{1-b}\le\epsilon
\]
uniformly in $a$.
Pack whole groups by next fit into dyadic supergroups of
size $H_0=2^{Th+o(h)}$, at least twice the largest group size after slot
duplication. This uses $O(n/H_0+1)$ supergroups. Take $K$ a sufficiently
small fixed dyadic fraction of $H_0$. List transports in each supergroup
in canonical order in batches of at most $K/2$, filling all but the last.
Reserving $K$ positions per nonempty batch uses in total at most
$2c_*n+O((n/H_0+1)K)\le3c_*n$ positions. This choice uses the uniform
lower bound on $c_*$ and $H_0=o(n)$.

Hardwire candidate batches for every supergroup, allowing all its slots
to be active. For each batch, Lemma~\ref{lem:slack-compaction} assigns active source and destination slots separately to two $K$-slot
lists by bounded-degree switches. Apply the same lemma to the set of
nonempty batches, assigning each entire pair of lists to one of
$\lceil6c_*n/K\rceil$ physical modules. The number of candidates is
$O(n/H_0+1)$, a bounded multiple of the number of modules, and at most
half the modules are needed for any summand. One
$\operatorname{Tag}(\log K,\lceil Yw\rceil)$ in each occupied module
has its tags set to the actual matching of the two lists. For the remaining
slots, set the tags to match the vacant destinations with vacant sources
in increasing order and fix all unused payload bits to zero. Idle modules
use the identity permutation and zero data. Thus every original link is
enforced with its prescribed multiplicity. There are between $6c_*n$
and $7c_*n$ installed slots for large $h$. This gives the weighted slot bound.

All assignment graphs are fixed advice of linear word size for fixed
parameters. Scans compute the active lists, the deterministic assignments
and the completed permutations in $O(n\,\mathrm{polylog}\,h)$ word
operations, including a scan of each candidate link to a $K$-slot
batch. Whole-batch switches transport $O(K\lceil Yw\rceil)$ bits per
incidence, while individual slot switches transport $O(\lceil Yw\rceil)$
bits. The number of incidences per position is bounded for fixed
parameters. Variable subfield shifts cost $O(w\log(2w))$ per site.
Together with the site networks, pins, selector copies and outside
catalyst ports these fit $O(nw\,\mathrm{polylog}\,h)$. For protocol
control, guess a timed RAM history and check each bounded-word step.
Sort the memory records, including initialization, by address and time
and scan them to verify that reads agree with the latest value. Pad the
history canonically and resolve ties by the step's fixed access order.
This leaves one history for each deterministic computation. Tag validity
can be tested by one more sort. A bounded
number of stable sorts use $O(n\,\mathrm{polylog}\,h)$ slots with
$O(h)$-bit records. Pad lengths to powers of two of address length at
least $h$, and split records into a bounded number of address-width
slices with shared tags. This gives the claimed control calls and
$O(nh\,\mathrm{polylog}\,h)$ further size.
\end{proof}

\subsection{Closing the induction}

Fix $u$ as above and choose $A$ sufficiently large in the interior of
its constant region. We prove the size and gap estimates together by
strong induction on the address length $h$. For every integer $h\ge2$
and every integer $w$ with
$h\le w\le h^A$, there are approximations with
\begin{equation}\label{eq:tag-uniform-width}
  \text{size}\le D_{\rm cost}\,2^h w h^{u(\log_h w)},\qquad
  \text{gap}\ge D_h=2^h\,2^{-C_{\rm gap}L(h)^2}.
\end{equation}
Choose $r_*<r'<u(1)$ and $\kappa>0$ small enough that $r'<u(a)$ on
$[1,1+\kappa]$. Fix $C_{\rm gap}>3$ strictly larger than the constant
in the width-$h$ estimate for $r'$. On this narrow-width interval,
slicing the payload into $\lceil w/h\rceil$ padded width-$h$ tests
proves \eqref{eq:tag-uniform-width} for sufficiently large $h$ with
any $D_{\rm cost}\ge1$. The strict margins absorb the slicing constants
and the polynomial absolute gap loss.

For $1+\kappa\le a\le A$, apply Lemma~\ref{lem:tag-shorter} with
$C_{\rm amp}\epsilon<1/4$. Use the $r'$ estimate for the control calls, whose payload
width equals their address length. Their gaps are $\omega(D_h)$ because their address
length is $h+O(\log h)$ and the exponent constant has a strict margin.
Every shorter type has $h'<h$ and $w'>h'$ for large $h$. If
$w'>(h')^A$, split into
$\lceil w'/\lfloor(h')^A\rfloor\rceil$ nearly equal integer slices.
There are boundedly many, each of width exponent $A+o(1)$ at height
$h'$, hence within the constant region. Treat each type as a single
batch of $s=S_{T,Y}/2^{h'}$ identical calls, separately for each
slice. Induction at $h'$ and Lemma~\ref{lem:tag-amplification} give
each batch gap at least
\[
          \frac{sD_{h'}}{2h'}
           =\Omega\left(\frac{n\,2^{-C_{\rm gap}L(h')^2}}{h'}\right)
           =\omega(D_h).
\]
Indeed $h'=Th+o(h)$ uniformly on the fixed menu, $T\le1/2$, and
$L(h)^2-L(h')^2\ge(2-o(1))\log h$. Also $\log s=O(h)$ and $D_{h'}$
exceeds the amplification threshold. The outer pruned gap is $\Omega(nw)$. Each amplified batch enters
the parent contraction as a single product gate. Thus
Lemma~\ref{lem:substitution} charges once per type and slice, and once
per control call. Their total number is bounded, so the resulting gap
is at least $D_h$ for sufficiently large $h$.

The ratio of the shorter per-slot size allowance to the parent allowance,
including the amplification factor, is
\[
 C_{\rm amp}\frac{w'(h')^{u(\log_{h'} w')}}{w h^{u(a)}}
        =(1+o(1)) C_{\rm amp}T^vY^{1-b}.
\]
In fact $\log h'=\log h+\log T+o(1)$ and
\[
 \log_{h'}w'-a=\frac{\log Y-a\log T+o(1)}{\log h}.
\]
Taylor expansion with $u'(a)=-b$ proves the formula uniformly on the
compact interval and fixed menu. For the split case the exponent is
constant and the slice widths sum to $w'$, giving the same allowance.
The shorter-call work consumes at most half the proposed size for large
$h$. For the remaining ordinary work and control costs we have
\[
 nw\,\mathrm{polylog}\,h+nh^{1+r'}\mathrm{polylog}\,h
                      =o(nwh^{u(a)}),
\]
since $u$ has positive minimum on the interval and
$a+u(a)\ge1+u(1)>1+r'$. This also pays for extra tag copying when
splitting calls. All thresholds are independent of $D_{\rm cost}\ge1$.
An enlargement of $D_{\rm cost}$ covers the finitely many smaller
formats using exact networks. This completes the induction.

\begin{proof}[Proof of Theorem~\ref{thm:cheaptag}]
Taking $A$ arbitrarily far into the constant region in
\eqref{eq:tag-uniform-width} gives $r(a)\le u(a)$ for every fixed $a$.
If $r_*>0$, we can choose $u$ to decrease fast enough just above
$a=1$, and to level off low enough at large $a$, that both terms in
the packing bound improve on $r_*$. To do so, choose $1-r_*<b_0<1$, and choose $t>0$ sufficiently
small and $0<e_0<(r_*+b_0-1)t$. On $[1,1+t]$ take
$u(a)=r_*+e_0-b_0(a-1)$ and extend it smoothly to the positive plateau
$r_*/3$ with $0\le-u'<1$. For small $t$, the linear piece stays above
the plateau and $t/(1+t)+r_*/3<r_*$. Also
$(t+u(1+t))/(1+t)<r_*$ by the choice of $e_0$. In
Proposition~\ref{prop:tag-width-packing} with $a=1+t$ substitute
$r(1+t)\le u(1+t)$ and $r_\infty\le r_*/3$. The right-hand side is
strictly less than $r_*$. Hence $r_*=0$, which gives the assertion
for every fixed $r>0$ by the definition of the profile.
\end{proof}
